\section{Discussion}
\subsection{What \sys{} Measures}
\sys{} evaluates whether agents can produce auditable evidence aligned with public vulnerability facts and real firmware artifacts, without assuming a prebuilt executable oracle for each CVE.

\subsection{Implications for Agent Evaluation}
Prepared environments hide a major class of agent failures. Future benchmarks should record traces and workspaces, require artifact provenance, distinguish evidence insufficiency from agent error, and reject substitutions that do not preserve task validity.

\subsection{Implications for Agent Design}
The results reveal an architecture-dependent ceiling: same-architecture devices (ARM64) allow \texttt{chroot}-based rehosting, while cross-architecture devices (MIPS, ARM32) requiring QEMU almost universally fail at R5. Future agent designs should prioritize cross-architecture emulation configuration---NVRAM initialization, library path resolution, and network bridge setup. Best-of-three scoring reveals high run-to-run variance (approximately +10 points over averages), with stronger models exhibiting larger variance---suggesting that consistency of peak performance is itself a capability dimension.

\subsection{Limitations and Ethics}
The dataset is balanced across three vulnerability classes but does not mirror the natural IoT distribution. The scoring skill is not a deterministic exploit oracle; human analysis validates evidence consistency. Ethically, \sys{} targets defensive reproduction using only publicly disclosed, remediated CVEs, conducted in isolated workspaces.
